% GENERATED FILE. Edit canonical lessons, then run npm run generate:books.
% canonical-source-hash: fnv1a64:f052ba1558493c53
% canonical-lessons: JA-C504-bungaku, JA-C504-shijin, JA-C504-garou, JA-C504-geijutsuka, JA-C504-ehon

\chapter{Books and Art}
\label{ch:ja-books-and-art}
\hlchaptermodality{504}

\begin{chapteropening}
\textbf{By the end of this chapter:} \emph{I can talk about literature and the people who make art.}

Complete the last lesson of chapter 504: I can talk about literature and the people who make art.
\end{chapteropening}

\section[\texorpdfstring{\ja{ぶんがく} (bungaku)}{ぶんがく (bungaku)}]{\ja{ぶんがく} (bungaku) — literature}

\label{lesson:JA-C504-bungaku}



\begin{quote}
Before the new one: say the Japanese for a musician, then the Japanese for to flutter.
\end{quote}

\subsection*{You'll want to know: \ja{ぶんがく}}
\textbf{\ja{ぶんがく}} — \emph{bungaku} — \textquotedblleft{}literature\textquotedblright{}.

\begin{quote}
\textbf{\ja{あには} \ja{にほんの} \ja{ぶんがくを} \ja{べんきょうしています。}} — \emph{ani wa nihon no bungaku o benkyō shiteimasu} — \textquotedblleft{}My elder brother studies Japanese literature.\textquotedblright{}
\end{quote}

\begin{cousinweb}
\textbf{\ja{ぶんがく}} (\emph{bungaku}) is Sino-Japanese (kango): two Chinese characters read with their on'yomi, readings borrowed from Chinese. They are \emph{bun}, \textquotedblleft{}sentence\textquotedblright{} (Middle Chinese \emph{mjun}) and \emph{gaku}, \textquotedblleft{}study\textquotedblright{} (\emph{haewk}) (KANJIDIC2; Middle Chinese in Baxter's transcription).
\end{cousinweb}

\begin{grammarlens}[title={What you've built}]
One more word for describing what you lived and felt.
\end{grammarlens}

\subsection*{Your turn}
\begin{itemize}
\raggedright
  \item \emph{Say it:} \emph{bungaku}
  \item \emph{Say it:} \emph{bungaku}, once more
  \item \emph{Say it:} \emph{ani wa nihon no bungaku o benkyō shiteimasu} — My elder brother studies Japanese literature.
  \item \emph{Recall:} say the Japanese for a musician, then the Japanese for to flutter, then say \emph{bungaku} again
\end{itemize}

\begin{tcolorbox}[breakable,colback=teal!4,colframe=teal!35!black,title={Before you move on}]
What does \ja{ぶんがく} mean? (\textquotedblleft{}Literature\textquotedblright{}.) Say it once more.
\end{tcolorbox}
\section[\texorpdfstring{\ja{しじん} (shijin)}{しじん (shijin)}]{\ja{しじん} (shijin) — a poet}

\label{lesson:JA-C504-shijin}



\begin{quote}
Before the new one: say the Japanese for to flutter, then the Japanese for literature.
\end{quote}

\subsection*{You'll want to know: \ja{しじん}}
\textbf{\ja{しじん}} — \emph{shijin} — \textquotedblleft{}a poet\textquotedblright{}.

\begin{quote}
\textbf{\ja{この} \ja{しじんは} \ja{ゆうめいです。}} — \emph{kono shijin wa yūmei desu} — \textquotedblleft{}This poet is famous.\textquotedblright{}
\end{quote}

\begin{cousinweb}
\textbf{\ja{しじん}} (\emph{shijin}) is Sino-Japanese (kango): two Chinese characters read with their on'yomi, readings borrowed from Chinese. They are \emph{shi}, \textquotedblleft{}poem\textquotedblright{} (Middle Chinese \emph{syi}) and \emph{jin}, \textquotedblleft{}person\textquotedblright{} (\emph{nyin}) (KANJIDIC2; Middle Chinese in Baxter's transcription).
\end{cousinweb}

\begin{grammarlens}[title={What you've built}]
One more word for describing what you lived and felt.
\end{grammarlens}

\subsection*{Your turn}
\begin{itemize}
\raggedright
  \item \emph{Say it:} \emph{shijin}
  \item \emph{Say it:} \emph{shijin}, once more
  \item \emph{Say it:} \emph{kono shijin wa yūmei desu} — This poet is famous.
  \item \emph{Recall:} say the Japanese for to flutter, then the Japanese for literature, then say \emph{shijin} again
\end{itemize}

\begin{tcolorbox}[breakable,colback=teal!4,colframe=teal!35!black,title={Before you move on}]
What does \ja{しじん} mean? (\textquotedblleft{}A poet\textquotedblright{}.) Say it once more.
\end{tcolorbox}
\section[\texorpdfstring{\ja{がろう} (garō)}{がろう (garō)}]{\ja{がろう} (garō) — an art gallery}

\label{lesson:JA-C504-garou}



\begin{quote}
Before the new one: say the Japanese for literature, then the Japanese for a poet.
\end{quote}

\subsection*{You'll want to know: \ja{がろう}}
\textbf{\ja{がろう}} — \emph{garō} — \textquotedblleft{}an art gallery\textquotedblright{}.

\begin{quote}
\textbf{\ja{まちの} \ja{がろうに} \ja{いきました。}} — \emph{machi no garō ni ikimashita} — \textquotedblleft{}I went to a gallery in town.\textquotedblright{}
\end{quote}

\begin{cousinweb}
\textbf{\ja{がろう}} (\emph{garō}) is Sino-Japanese (kango): two Chinese characters read with their on'yomi, readings borrowed from Chinese. They are \emph{ga}, \textquotedblleft{}picture\textquotedblright{} (Middle Chinese \emph{hweaH}) and \emph{rō}, \textquotedblleft{}corridor\textquotedblright{} (\emph{lang}) (KANJIDIC2; Middle Chinese in Baxter's transcription).
\end{cousinweb}

\begin{grammarlens}[title={What you've built}]
One more word for describing what you lived and felt.
\end{grammarlens}

\subsection*{Your turn}
\begin{itemize}
\raggedright
  \item \emph{Say it:} \emph{garō}
  \item \emph{Say it:} \emph{garō}, once more
  \item \emph{Say it:} \emph{machi no garō ni ikimashita} — I went to a gallery in town.
  \item \emph{Recall:} say the Japanese for literature, then the Japanese for a poet, then say \emph{garō} again
\end{itemize}

\begin{tcolorbox}[breakable,colback=teal!4,colframe=teal!35!black,title={Before you move on}]
What does \ja{がろう} mean? (\textquotedblleft{}An art gallery\textquotedblright{}.) Say it once more.
\end{tcolorbox}
\section[\texorpdfstring{\ja{げいじゅつか} (geijutsuka)}{げいじゅつか (geijutsuka)}]{\ja{げいじゅつか} (geijutsuka) — an artist}

\label{lesson:JA-C504-geijutsuka}



\begin{quote}
Before the new one: say the Japanese for a poet, then the Japanese for an art gallery.
\end{quote}

\subsection*{You'll want to know: \ja{げいじゅつか}}
\textbf{\ja{げいじゅつか}} — \emph{geijutsuka} — \textquotedblleft{}an artist\textquotedblright{}.

\begin{quote}
\textbf{\ja{かのじょは} \ja{げいじゅつかです。}} — \emph{kanojo wa geijutsuka desu} — \textquotedblleft{}She is an artist.\textquotedblright{}
\end{quote}

\begin{cousinweb}
\textbf{\ja{げいじゅつか}} (\emph{geijutsuka}) is Sino-Japanese (kango): three Chinese characters read with their on'yomi, readings borrowed from Chinese. They are \emph{gei}, \textquotedblleft{}art\textquotedblright{} (Middle Chinese \emph{ngjiejH}), \emph{jutsu}, \textquotedblleft{}technique\textquotedblright{} (\emph{zywit}) and \emph{ka}, \textquotedblleft{}expert\textquotedblright{} (\emph{kae}) (KANJIDIC2; Middle Chinese in Baxter's transcription).
\end{cousinweb}

\begin{grammarlens}[title={What you've built}]
One more word for describing what you lived and felt.
\end{grammarlens}

\subsection*{Your turn}
\begin{itemize}
\raggedright
  \item \emph{Say it:} \emph{geijutsuka}
  \item \emph{Say it:} \emph{geijutsuka}, once more
  \item \emph{Say it:} \emph{kanojo wa geijutsuka desu} — She is an artist.
  \item \emph{Recall:} say the Japanese for a poet, then the Japanese for an art gallery, then say \emph{geijutsuka} again
\end{itemize}

\begin{tcolorbox}[breakable,colback=teal!4,colframe=teal!35!black,title={Before you move on}]
What does \ja{げいじゅつか} mean? (\textquotedblleft{}An artist\textquotedblright{}.) Say it once more.
\end{tcolorbox}
\section[\texorpdfstring{\ja{えほん} (ehon)}{えほん (ehon)}]{\ja{えほん} (ehon) — a picture book}

\label{lesson:JA-C504-ehon}



\begin{quote}
Before the new one: say the Japanese for an art gallery, then the Japanese for an artist.
\end{quote}

\subsection*{You'll want to know: \ja{えほん}}
\textbf{\ja{えほん}} — \emph{ehon} — \textquotedblleft{}a picture book\textquotedblright{}.

\begin{quote}
\textbf{\ja{ははは} \ja{こどもに} \ja{えほんを} \ja{よみます。}} — \emph{haha wa kodomo ni ehon o yomimasu} — \textquotedblleft{}My mother reads picture books to the children.\textquotedblright{}
\end{quote}

\begin{cousinweb}
\textbf{\ja{えほん}} (\emph{ehon}) is Sino-Japanese (kango): two Chinese characters read with their on'yomi, readings borrowed from Chinese. They are \emph{e}, \textquotedblleft{}picture\textquotedblright{} and \emph{hon}, \textquotedblleft{}book\textquotedblright{} (Middle Chinese \emph{pwonX}) (KANJIDIC2; Middle Chinese in Baxter's transcription).
\end{cousinweb}

\begin{grammarlens}[title={What you've built}]
One more word for describing what you lived and felt.
\end{grammarlens}

\subsection*{Your turn}
\begin{itemize}
\raggedright
  \item \emph{Say it:} \emph{ehon}
  \item \emph{Say it:} \emph{ehon}, once more
  \item \emph{Say it:} \emph{haha wa kodomo ni ehon o yomimasu} — My mother reads picture books to the children.
  \item \emph{Recall:} say the Japanese for an art gallery, then the Japanese for an artist, then say \emph{ehon} again
\end{itemize}

\begin{tcolorbox}[breakable,colback=teal!4,colframe=teal!35!black,title={Before you move on}]
What does \ja{えほん} mean? (\textquotedblleft{}A picture book\textquotedblright{}.) Say it once more.
\end{tcolorbox}
